\textbf{\textit{Lambda expressions}} are composed of:

\begin{itemize}
\tightlist
\item
  variables~\iMbox{v_1,v_2, \dots};
\item
  the abstraction symbols \iMbox{\lambda} (lambda) and \iMbox{.} (dot);
\item
  parentheses \iMbox{()}. The set of lambda
  expressions,~\iMbox{\Lambda}, can be~\ul{defined inductively}:
\item
  Variable: \iMbox{x \ \text{is a variable} \implies x\in \Lambda};
\item
  Abstraction:
  \iMbox{\displaystyle x \ \text{is a variable and} \ M\in \Lambda \implies ( \lambda x . M) \in \Lambda};
\item
  Application:
  \iMbox{\displaystyle M,N \in \Lambda \implies ( M \ N ) \in \Lambda};
  Or in~\ul{Backus--Naur form}, written as
\end{itemize}

\begin{lstlisting}[language=bnf, linewidth=100pt]
 <expression>  ::= <abstraction>
                 | <application> 
                 | <variable>
 <abstraction> ::= λ <variable> . <expression>
 <application> ::= ( <expression> <expression> )
 <variable>    ::= v1 | v2 | ...
\end{lstlisting}

\section*{Conventional Notation}

To keep the notation of lambda expressions uncluttered, the following
conventions are usually applied:

\begin{itemize}
\tightlist
\item
  Outermost parentheses are dropped:~\iMbox{M\ N}~instead of
  \iMbox{(M\ N)}
\item
  Applications are assumed to be \ul{left
  associative}:~\iMbox{M\ N \ P= ((M \ N) \ P)}
\item
  When all variables are single-letter, the space in applications may be
  omitted:~\iMbox{xyz = x \ y \ z}
\item
  A \ul{lambda abstraction} has a \ul{lower precedence} than an
  \ul{application}: \iMbox{\lambda x.M \ N = \lambda x. (M \ N)}
\item
  A \ul{sequence of abstractions} is \ul{contracted}:
  \iMbox{\lambda x. \lambda y. \lambda z. M = \lambda xyz. M}
\end{itemize}

\section*{Free and bound variables}

The abstraction operator \iMbox{\lambda} \ul{binds its variable}
wherever it occurs \ul{in the body} of the abstraction. Variables that
fall \ul{within the scope} of an abstraction are said to
be~\textbf{\textit{bound}}. All other variables are
called~\textbf{\textit{free}}. A variable is bound by its ``nearest''
abstraction.

\begin{itemize}
\tightlist
\item
  e.g.~in \iMbox{\lambda y. x x y}, the \iMbox{y} is
  \textbf{\textit{bound}} and the \iMbox{x} is \textbf{\textit{free}}
\item
  e.g.~in \iMbox{\lambda x. y (\lambda x. xz)} the \iMbox{x} is
  \textbf{\textit{bound}} to the second lambda
\end{itemize}

\subsection*{\texorpdfstring{Set of free variables
\iMbox{\text{FV}(M)}}{Set of free variables }}

The \textbf{\textit{set of free variables}} of a \ul{lambda expression}
\iMbox{M} is denoted as \iMbox{\text{FV}(M)} and \ul{defined by
recursion} on the structure of the terms, as follows:

\begin{enumerate}
\def\labelenumi{\arabic{enumi}.}
\tightlist
\item
  \iMbox{\text{FV}(x) = \{ x \}}
\item
  \iMbox{\text{FV}(\lambda x. M) = \text{FV}(M) \setminus \{ x \}}
\item
  \iMbox{\text{FV}(M \ N) = \text{FV}(M) \cup \text{FV}(N)}
\end{enumerate}

\subsection*{Closed lambda expressions, i.e.~combinators}

An expression that contains \ul{no free variables} is said to
be~\textbf{\textit{closed}}. They are also known as
\textbf{\textit{combinators}} and are equivalent to terms
in~\ul{combinatory logic}.

\section*{Reduction}

The meaning of \ul{lambda expressions} is defined by how expressions can
be \textbf{\textit{reduced}}.

There are three kinds of \textbf{\textit{reduction}}:

\begin{itemize}
\tightlist
\item
  \ul{\iMbox{\alpha}-conversion}: changing bound variables
  (\textbf{alpha});
\item
  \ul{\iMbox{\beta}-reduction}: applying functions to their arguments
  (\textbf{beta});
\item
  \ul{\iMbox{\eta}-reduction}: which captures a notion of extensionality
  (\textbf{eta}).
\end{itemize}

\subsection*{Equivalences}

We also speak of the resulting \textbf{\textit{equivalences}}: two
expressions are~\textbf{\textit{equivalent}}, if they can be
\textbf{\textit{reduced}} into the same expression

\begin{itemize}
\tightlist
\item
  \iMbox{\displaystyle M_{1} =_{\alpha} M_{2} \iff \exists M'. \ M_{1} \to_{\alpha}^{*} M' \ \text{and} \ M_{2} \to_{\alpha}^{*} M'}
\item
  \iMbox{\displaystyle M_{1} =_{\beta} M_{2} \iff \exists M'. \ M_{1} \to_{\beta}^{*} M' \ \text{and} \ M_{2} \to_{\beta}^{*} M'}
\item
  \iMbox{\displaystyle M_{1} =_{\eta} M_{2} \iff \exists M'. \ M_{1} \to_{\eta}^{*} M' \ \text{and} \ M_{2} \to_{\eta}^{*} M'}
  And we can even combine them to achieve more \textit{complete} notions
  of \textbf{\textit{equivalence}}
\item
  \iMbox{\displaystyle M_{1} =_{\beta \eta} M_{2} \iff \exists M'. \ M_{1} \to_{\beta \eta}^{*} M' \ \text{and} \ M_{2} \to_{\beta \eta}^{*} M'}
\end{itemize}

\subsection*{\texorpdfstring{\iMbox{\alpha}-conversion}{-conversion}}

\textbf{\textit{Alpha-conversion}}, sometimes known as
\textbf{\textit{alpha-renaming}}, allows bound \ul{variable names} to be
changed. It is subject to these rules

\begin{itemize}
\tightlist
\item
  ~only rename variable occurrences bound by the \ul{same abstraction}

  \begin{itemize}
  \tightlist
  \item
    e.g.~\iMbox{\displaystyle \lambda x. \lambda x. x =_{\alpha} \lambda y . \lambda x. x}
    or
    \iMbox{\displaystyle \lambda x. \lambda x. x =_{\alpha} \lambda x. \lambda y. y}
    but \ul{not}
    \iMbox{\displaystyle \lambda x. \lambda x. x =_{\alpha} \lambda y . \lambda x. y}
  \end{itemize}
\item
  you \ul{cannot} rename a \ul{variable} if it will result in a
  \ul{different} \iMbox{\lambda} capturing it

  \begin{itemize}
  \tightlist
  \item
    e.g.~you \ul{cannot} do
    \iMbox{\lambda x. \lambda y. x =_{\alpha} \lambda y. \lambda y. y}
  \end{itemize}
\end{itemize}

\subsubsection*{Substitution}

\textbf{\textit{Substitution}}, written \iMbox{E[R/x]}, is the process
of replacing all \ul{free occurrences} of variable \iMbox{x} in the
expression \iMbox{E} with expression \iMbox{R}.
\textbf{\textit{Substitution}} on terms of the lambda calculus is
defined by recursion on the structure of terms, as follows (note:
\iMbox{x} and \iMbox{y} are \textit{only} \ul{variables} while \iMbox{M}
and \iMbox{N} are any \iMbox{\lambda}-expression) \[\displaystyle 
\begin{align}
x[\textcolor{red}{M}/y] &\equiv \begin{cases}
\textcolor{red}{M} & x=y \\
x & x \neq y
\end{cases} \\
(\textcolor{red}{M_{1}} \ \textcolor{green}{M_{2}})[\textcolor{blue}{N} / x] &\equiv (\textcolor{red}{M_{1}}[\textcolor{blue}{N} / x]) \ (\textcolor{green}{M_{2}} [\textcolor{blue}{N} / x]) \\
(\lambda x. \textcolor{red}{M})[\textcolor{green}{N} /y] &\equiv \begin{cases}
\lambda x. \textcolor{red}{M} & x=y \\
\lambda x. (\textcolor{red}{M}[\textcolor{green}{N} /y]) & x\neq y \ \text{and} \ x \not\in \text{FV}(N) \\
\lambda z. (\textcolor{red}{M}[z/x][\textcolor{green}{N} /y]) & \begin{aligned}
&x\neq y \ \text{and} \ x \in \text{FV}(N), \\
&\text{for some} \ z \neq x, z \neq y, z \not\in \text{FV}(M), z \not\in \text{FV}(N) 
\end{aligned}
\end{cases} \\
\end{align}
\] To \textbf{\textit{substitute}} into a lambda abstraction, it is
sometimes necessary to \ul{\iMbox{\alpha}-convert} the expression.

\begin{itemize}
\tightlist
\item
  e.g.~its \ul{wrong} for \iMbox{\displaystyle (\lambda x.y)[x / y]} to
  result in \iMbox{(\lambda x.x)}, because \iMbox{x} was meant to be
  \ul{free not bound}
\item
  instead the correct substitution is \iMbox{(\lambda z.x)} up to
  \ul{\iMbox{\alpha}-equivalence}
\end{itemize}

\subsection*{\texorpdfstring{\iMbox{\beta}-reduction}{-reduction}}

\textbf{\textit{Beta-reduction}} captures the idea of \ul{function
application}. It can be defined in terms of \ul{substitution} where the
\textbf{\textit{\iMbox{\beta}-reduction}} of
\iMbox{\displaystyle (\lambda x. M) \ N} is \iMbox{M[N / x]}.

\begin{itemize}
\tightlist
\item
  e.g.~\iMbox{\displaystyle ((\lambda n. n \times 2) \ 7) \longrightarrow _{\beta} 7 \times 2}
  From there, you can derive a couple more rules to make performing
  \textbf{\textit{\iMbox{\beta}-reduction}} more convenient in different
  scenarios: \[\large 
  \begin{gathered}
  \begin{prooftree} 
  \AxiomC{$\textcolor{red}{M} \longrightarrow_{\beta} \textcolor{green}{M'}$}
  \UnaryInfC{$\lambda x. \textcolor{red}{M} \longrightarrow_{\beta} \lambda x. \textcolor{green}{M'}$}  
  \end{prooftree} &
  \begin{prooftree} 
  \AxiomC{$\textcolor{red}{M} \longrightarrow_{\beta} \textcolor{green}{M'}$}
  \UnaryInfC{$\textcolor{red}{M}N \longrightarrow_{\beta} \textcolor{green}{M'}N$}  
  \end{prooftree} &
  \begin{prooftree} 
  \AxiomC{$\textcolor{red}{N} \longrightarrow_{\beta} \textcolor{green}{N'}$}
  \UnaryInfC{$M\textcolor{red}{N} \longrightarrow_{\beta} M\textcolor{green}{N'}$}  
  \end{prooftree} &
  \begin{prooftree}
  \AxiomC{$\textcolor{red}{M} =_{\alpha} \textcolor{red}{M'}$}
  \AxiomC{$\textcolor{red}{M'} \longrightarrow_{\beta} \textcolor{green}{N'}$}
  \AxiomC{$\textcolor{green}{N'} =_{\alpha} \textcolor{green}{N}$}
  \TrinaryInfC{$\textcolor{red}{M} \longrightarrow_{\beta} \textcolor{green}{N}$}  
  \end{prooftree}
  \end{gathered}
  \] \#\#\# Multiple \iMbox{\beta}-reductions:
  \iMbox{\displaystyle \to_{\beta}^{*}} If we can reduce expression
  \iMbox{M} into expression \iMbox{M} with \textbf{\textit{multiple
  \iMbox{\beta}-reductions}}, we express that fact with
  \iMbox{\displaystyle M \to_{\beta}^{*} M'}, with the \iMbox{*} being
  shorthand for the sequence of
  \textbf{\textit{\iMbox{\beta}-reductions}}.
\end{itemize}

\subsubsection*{\texorpdfstring{Reflexivity of
\ul{\iMbox{\alpha}-conversion}}{Reflexivity of -conversion}}

\begin{figure}
\centering
\textcolor{red}{\textbf{Pasted image 20241121132159.png}}
\caption{100}
\end{figure}

\subsubsection*{\texorpdfstring{Transitivity of
\ul{\iMbox{\beta}-reduction}}{Transitivity of -reduction}}

\begin{figure}
\centering
\textcolor{red}{\textbf{Pasted image 20241121132238.png}}
\caption{250}
\end{figure}

\subsection*{\texorpdfstring{\iMbox{\eta}-reduction}{-reduction}}

\textbf{\textit{Eta-reduction}} expresses the idea
of~\ul{extensionality}, which in this context is that two functions are
the \textbf{same} IFF~they give the \ul{same result} for \ul{all
arguments}. \[\large
\begin{prooftree} 
\AxiomC{$\displaystyle x \not\in \text{FV}(\textcolor{red}{M})$}
\UnaryInfC{$\displaystyle  (\lambda x. \textcolor{red}{M} x) \longrightarrow_{\eta} \textcolor{red}{M}$}  
\end{prooftree}
\]

\subsection*{Redexes and Reducts}

The term~\textbf{\textit{redex}}, short for~\textbf{\textit{reducible
expression}}, refers to subterms that can be reduced by one of the
\ul{reduction rules}. For example, \iMbox{(\lambda x.M)}~is a
\textbf{\textit{\iMbox{\beta}-redex}} in expressing the
\ul{substitution} of~\iMbox{N}~for~\iMbox{x}~in~\iMbox{M};
if~\iMbox{x}~is \ul{not free} in~\iMbox{M},~\iMbox{\lambda x.Mx}~is an
\textbf{\textit{\iMbox{\eta}-redex}}. The expression to which a
\ul{redex} reduces is called its \textbf{\textit{reduct}}; using the
previous example, the \textbf{\textit{reducts}} of these expressions are
respectively~\iMbox{M[N/x]}~and~\iMbox{M}.

\subsubsection*{\texorpdfstring{Innermost and Outermost
\iMbox{\beta}-redexes}{Innermost and Outermost -redexes}}

Take the \ul{redex} \iMbox{\displaystyle E = (\lambda x. M) \ N},

\begin{itemize}
\tightlist
\item
  any \ul{redex} in \iMbox{M} or \iMbox{N} is \textbf{\textit{inside}}
  of the \ul{redex} \iMbox{E}\\
\item
  the \ul{redex} \iMbox{E} is \textbf{\textit{outside}} of any
  \ul{redex} in \iMbox{M} or \iMbox{N} Then we can use these notions to
  say that
\item
  A \ul{redex} is \textbf{\textit{outermost}} if there are no
  \textit{larger} \ul{redexes} \textbf{\textit{outside}} it which
  \textit{contain} it
\item
  A \ul{redex} is \textbf{\textit{innermost}} if there are no
  \ul{redexes} \textbf{\textit{inside}} it which it \textit{contains}
  \textcolor{red}{\textbf{Pasted image 20241121142017.png}}
\end{itemize}

\subsubsection*{\texorpdfstring{Head \& Internal
\iMbox{\beta}-redexes}{Head \& Internal -redexes}}

A \ul{redex} \iMbox{r} is in the \textbf{\textit{head position}} in an
expression \iMbox{E} if it has the following shape:

\begin{itemize}
\tightlist
\item
  \iMbox{\displaystyle \lambda x_{0}\dots \lambda x_{n} . \textcolor{red}{(\lambda x. N) \ M} \ M_{0} \dots M_{m}},
  where \iMbox{n,m \geq 0}

  \begin{itemize}
  \tightlist
  \item
    Here,
    \iMbox{\displaystyle r \equiv \textcolor{red}{(\lambda x. N) \ M}}
    is the \textbf{\textit{head position}} and is called the
    \textbf{\textit{head \iMbox{\beta}-redex}}
  \end{itemize}
\item
  A \textbf{\textit{head \iMbox{\beta}-reduction}} is a
  \ul{\iMbox{\beta}-reduction} applied to the \textbf{\textit{head
  position}}, e.g.~here applied to
  \iMbox{\displaystyle r \equiv \textcolor{red}{(\lambda x. N) \ M}}

  \begin{itemize}
  \tightlist
  \item
    e.g.~here it would be
    \iMbox{\displaystyle \lambda x_{0}\dots \lambda x_{n} . \textcolor{red}{(\lambda x. N) \ M} \ M_{0} \dots M_{m} \ \longrightarrow \ \lambda x_{0}\dots \lambda x_{n} . \textcolor{red}{N[M / x]} \ M_{0} \dots M_{m}}
  \end{itemize}
\item
  Any other \ul{\iMbox{\beta}-reduction} is called an
  \textbf{\textit{internal \iMbox{\beta}-reduction}} and that \ul{redex}
  is called an \textbf{\textit{internal \iMbox{\beta}-redex}}
\end{itemize}

\section*{Confluence (Church--Rosser property)}

The property of \textbf{\textit{confluence}} is a general \ul{rewriting
systems} property. The idea of \textbf{\textit{confluence}} is that, if
\iMbox{a} \ul{reduces} to both \iMbox{b} and \iMbox{c} in \iMbox{\geq 0}
rewrite steps then both \ul{reducts} must in turn reduce to some
common~\iMbox{d}.
\textcolor{red}{\textbf{Pasted image 20241121132641.png}}\textcolor{red}{\textbf{Pasted image 20241121133521.png}}
When applied to \ul{lambda calculus reductions}, this property is known
as the \textbf{\textit{Church--Rosser property}}. It has been stated and
proven for

\begin{itemize}
\tightlist
\item
  \ul{Confluence} of \ul{\iMbox{\beta}-reductions}:
  \iMbox{\displaystyle \forall M,M_{1},M_{2}. \ \ M \to_{\beta}^{*} M_{1} \ \text{and} \ M \to_{\beta}^{*} M_{2} \implies \exists M'. \ M_{1} \to_{\beta}^{*} M' \ \text{and} \ M_{2} \to_{\beta}^{*} M'}
\item
  \ul{Confluence} of \ul{\iMbox{\eta}-reductions}:
  \iMbox{\displaystyle \forall M,M_{1},M_{2}. \ \ M \to_{\eta}^{*} M_{1} \ \text{and} \ M \to_{\eta}^{*} M_{2} \implies \exists M'. \ M_{1} \to_{\eta}^{*} M' \ \text{and} \ M_{2} \to_{\eta}^{*} M'}
\item
  \ul{Confluence} of \iMbox{\beta \eta}\ul{-reductions}:
  \iMbox{\displaystyle \forall M,M_{1},M_{2}. \ \ M \to_{\beta \eta}^{*} M_{1} \ \text{and} \ M \to_{\beta \eta}^{*} M_{2} \implies \exists M'. \ M_{1} \to_{\beta \eta}^{*} M' \ \text{and} \ M_{2} \to_{\beta \eta}^{*} M'}
\end{itemize}

\section*{Normalization}

A lambda expression that \ul{cannot} be \ul{reduced further} is in
\textbf{\textit{normal form}}. You can still apply
\ul{\iMbox{\alpha}-conversion} to expressions in \textbf{\textit{normal
form}}, in fact all \textbf{\textit{normal forms}} that can be
\textit{converted} into each other by \ul{\iMbox{\alpha}-conversion} are
defined to be \ul{equal}. Here is a summary of \textbf{\textit{normal
forms}}:

\begin{longtable}[]{@{}
  >{\raggedright\arraybackslash}p{(\columnwidth - 2\tabcolsep) * \real{0.3544}}
  >{\raggedright\arraybackslash}p{(\columnwidth - 2\tabcolsep) * \real{0.6456}}@{}}
\toprule\noalign{}
\begin{minipage}[b]{\linewidth}\raggedright
Normal Form Type
\end{minipage} & \begin{minipage}[b]{\linewidth}\raggedright
Definition
\end{minipage} \\
\midrule\noalign{}
\endhead
\bottomrule\noalign{}
\endlastfoot
\ul{\iMbox{\beta \eta}-Normal Form} \textbar{} No
\ul{\iMbox{\beta}-reductions} or \ul{\iMbox{\eta}-reductions} are
possible \textbar{} \textbar{} \ul{\iMbox{\beta}-Normal Form} \textbar{}
No \ul{\iMbox{\beta}-reductions} are possible & \\
\ul{Head Normal Form} & In the form of a lambda abstraction \ul{whose
body is not reducible} \\
\ul{Weak Head Normal Form} & In the form of a lambda abstraction \\
\end{longtable}

\^{}these \textbf{\textit{normal forms}} form a ``subset'' hierarchy:
\iMbox{\beta \eta \text{-NF} \subset \beta \text{-NF} \subset \text{HNF} \subset \text{WHNF}}

Normal forms \textit{(both \ul{\iMbox{\beta \eta}-NFs} and
\ul{\iMbox{\beta}-NFs})} are \ul{unique/equivalent} up to an
\ul{\iMbox{\alpha}-conversion},

\begin{itemize}
\tightlist
\item
  \iMbox{\forall M,N_{1},N_{2}. \ \ M \to_{\beta \eta}^{*} N_{1} \ \text{and} \ M \to_{\beta \eta}^{*} N_{2} \ \text{and} \ N_{1},N_{2} \ \text{are in} \ \beta \eta \text{-normal form} \implies N_{1} =_{\alpha} N_{2}}
\item
  \iMbox{\forall M,N_{1},N_{2}. \ \ M \to_{\beta}^{*} N_{1} \ \text{and} \ M \to_{\beta}^{*} N_{2} \ \text{and} \ N_{1},N_{2} \ \text{are in} \ \beta \text{-normal form} \implies N_{1} =_{\alpha} N_{2}}
  This is \textit{merely} a corollary to them \ul{both} having the
  \ul{Church-Rosser property},
\item
  If expression \iMbox{M} has two \textbf{\textit{NFs}}
  \iMbox{N_{1},N_{2}} then \iMbox{M} \ul{must} have produced them via
  \ul{reductions} \textit{(\iMbox{\beta \eta} or \iMbox{\beta}
  respectively)}
\item
  Therefore by \ul{confluence}, \iMbox{N_{1}} and \iMbox{N_{2}} must
  \ul{both} arrive at some \ul{common} \iMbox{M'} with \iMbox{\geq 0}
  \ul{reductions} *(\iMbox{\beta \eta} or \iMbox{\beta} respectively)
\item
  But since \iMbox{N_{1}}, \iMbox{N_{2}} are \textbf{\textit{NFs}} they
  contain no \ul{redexes} \textit{(\iMbox{\beta \eta} or \iMbox{\beta}
  respectively)} so \iMbox{N_{1}}, \iMbox{N_{2}} arrive at \iMbox{M'}
  with exactly \iMbox{0} \ul{reductions} \textit{(\iMbox{\beta \eta} or
  \iMbox{\beta} respectively)}
\item
  Meaning they were either \ul{already equal}
  \textit{(i.e.~\iMbox{N_{1} \equiv M \equiv N_{2}})} or they needed to
  \ul{\iMbox{\alpha}-convert}
  \textit{(i.e.~\iMbox{N_{1} =_{a} M' =_{a} N_{2}})} making them
  \ul{equivalent} up to an \ul{\iMbox{\alpha}-conversion}
  \textbf{\textit{NOTE:}} not \textit{all} expressions
  \textit{necessarily} have a \textbf{\textit{normal form}},
  e.g.~\iMbox{(\lambda x. xx)(\lambda x.xx) \longrightarrow_{\beta} (\lambda x. xx)(\lambda x.xx)}
  so this expression will \textit{always} have a
  \ul{\iMbox{\beta}-redex}
\end{itemize}

\subsection*{\texorpdfstring{\iMbox{\beta \eta}-Normal
Form}{-Normal Form}}

A term is in~\textbf{\textit{\iMbox{\beta \eta}-Normal Form}}~if no
\ul{\iMbox{\beta}-reductions} or \ul{\iMbox{\eta}-reductions} are
possible, or alternatively, if it contains no \ul{\iMbox{\beta}-redexes
or \iMbox{\eta}-redexes}.

\begin{itemize}
\tightlist
\item
  \iMbox{M \ \text{is in} \ \beta \eta \text{-normal form} \ \ \coloneqq \ \ \neg \exists M'. \ M \to_{\beta} M' \ \text{or} \ M \to_{\eta} M'}
\item
  \iMbox{M \ \text{has a} \ \beta \eta \text{-normal form} \ \ \coloneqq \ \ \exists M'. \ M \to_{\beta \eta}^{*} M' \ \text{and} \ M' \ \text{is in} \ \beta \eta \text{-normal form}}
\end{itemize}

\subsection*{\texorpdfstring{\iMbox{\beta}-Normal Form}{-Normal Form}}

A term is in~\textbf{\textit{\iMbox{\beta}-Normal Form}}~if no
\ul{\iMbox{\beta}-reductions} are possible, or alternatively, if it
contains no \ul{\iMbox{\beta}-redexes}.

\begin{itemize}
\tightlist
\item
  \iMbox{M \ \text{is in} \ \beta \text{-normal form} \ \ \coloneqq \ \ \neg \exists M'. \ M \to_{\beta} M'}
\item
  \iMbox{M \ \text{has a} \ \beta \text{-normal form} \ \ \coloneqq \ \ \exists M'. \ M \to_{\beta}^{*} M' \ \text{and} \ M' \ \text{is in} \ \beta \text{-normal form}}
\end{itemize}

\subsection*{Head Normal Form (HNF)}

A term is in \textbf{\textit{Head Normal Form (HNF)}}~if it contains no
\ul{head \iMbox{\beta}-reducts}.

\begin{itemize}
\tightlist
\item
  Terms in \textbf{\textit{Head Normal Form (HNF)}} therefore have the
  following shape

  \begin{itemize}
  \tightlist
  \item
    \iMbox{\displaystyle \lambda x_{0}\dots \lambda x_{n} . \textcolor{red}{x} \ M_{0} \dots M_{m}},
    where \iMbox{x} is a \ul{variable} and \iMbox{n,m \geq 0}
  \end{itemize}
\item
  \textbf{\textit{Head Normal Form (HNF)}} terms are not
  \textit{necessarily} also \ul{\iMbox{\beta \eta}-NF} or
  \ul{\iMbox{\beta}-NF}
\end{itemize}

\subsection*{Weak Head Normal Form (WHNF)}

A term is in \textbf{\textit{Weak Head Normal Form (WHNF)}}~if its
either

\begin{itemize}
\tightlist
\item
  in \ul{head normal form}, or
\item
  it is a \ul{lambda abstraction} This means a WHNF term can contain
  \ul{\iMbox{\beta}-redexes} in the body of the lambda.
\end{itemize}

\section*{Reduction Strategies}

\textcolor{red}{\textbf{Pasted image 20241121182543.png}} \^{}quick
summary of main strategies

\subsection*{Normal Order}

\begin{itemize}
\tightlist
\item
  \ul{\iMbox{\beta}-reduces} the \ul{leftmost + outermost redex} first

  \begin{itemize}
  \tightlist
  \item
    i.e.~whenever possible, arguments are substituted into the body of
    an abstraction before the arguments are reduced
  \end{itemize}
\item
  \ul{Always} reduces a term to its \ul{\iMbox{\beta}-normal form}
  \textit{(if it exists)}
\item
  Can perform computations in unevaluated function bodies
\item
  Is not used by any programming language
\end{itemize}

\subsection*{Applicative Order}

\begin{itemize}
\tightlist
\item
  \ul{\iMbox{\beta}-reduces} the \ul{leftmost + innermost redex} first

  \begin{itemize}
  \tightlist
  \item
    therefore, ~a function's arguments are always reduced before they
    are substituted into the function
  \end{itemize}
\item
  Does not always reduce a term to its \ul{\iMbox{\beta}-normal form}
  \textit{(even if exists)}

  \begin{itemize}
  \tightlist
  \item
    e.g.~\iMbox{\displaystyle ( \ \ \ \lambda x. y \ \ (\lambda z. (zz) \lambda z. (zz) \ \ \ )}
    is reduced to itself by \textbf{\textit{applicative order}},
  \item
    while \ul{normal order} reduces it to its \ul{\iMbox{\beta}-normal
    form} of \iMbox{y}
  \end{itemize}
\end{itemize}

\subsection*{\texorpdfstring{Full
\ul{\iMbox{\beta}-reductions}}{Full -reductions}}

\begin{itemize}
\tightlist
\item
  Any \ul{\iMbox{\beta}-redex} can be \ul{\iMbox{\beta}-reduced} at any
  time
\item
  This means essentially the lack of any \textit{particular}
  \ul{reduction strategy}

  \begin{itemize}
  \tightlist
  \item
    With regard to \ul{reducibility}, ``all bets are off''
  \end{itemize}
\end{itemize}

\subsection*{Call by Name}

\begin{itemize}
\tightlist
\item
  Like \ul{normal order}, but \ul{no reductions are performed inside
  abstractions}

  \begin{itemize}
  \tightlist
  \item
    e.g.~\iMbox{\lambda x. (\lambda y.y) x} will \ul{not} be
    \ul{reduced} by this strategy
  \end{itemize}
\item
  Does not always reduce a term to its \ul{\iMbox{\beta}-normal form}
  \textit{(even if exists)}
\item
  Passes the function parameters unevaluated into the function body
\item
  Evaluates the passed function parameter on each use
\item
  Is used, with some variations, by, for example, Algol60, Haskell, R,
  and LaTeX
\end{itemize}

\subsection*{Call by Value}

\begin{itemize}
\tightlist
\item
  Like \ul{applicative order}, but \ul{no reductions are performed
  inside abstractions}

  \begin{itemize}
  \tightlist
  \item
    e.g.~\iMbox{\lambda x. (\lambda y.y) x} will \ul{not} be
    \ul{reduced} by this strategy
  \end{itemize}
\item
  Does not always reduce a term to its \ul{\iMbox{\beta}-normal form}
  \textit{(even if exists)}
\item
  Evaluates the function parameters before passing them into the
  function body
\item
  Terminates less often than \ul{call by name}, but evaluates parameters
  only once
\item
  Is used, with some variations, by, for example, C, Scheme, and OCamlX
\end{itemize}

\subsection*{Call by Need}

\begin{itemize}
\tightlist
\item
  As \ul{call by name}, but function applications that would duplicate
  terms instead name the argument

  \begin{itemize}
  \tightlist
  \item
    The argument may be evaluated ``when needed''
  \item
    at which point the name binding is updated with the reduced value
  \end{itemize}
\item
  This can save time compared to \ul{normal order} evaluation.
\end{itemize}

\subsection*{Optimal Reduction}

As \ul{normal order}, but \ul{computations} that have the \ul{same
label} are \ul{reduced simultaneously}
